\documentclass[10pt]{article}
\usepackage[utf8]{inputenc}
\usepackage[T1]{fontenc}
\usepackage[paperwidth=384pt,paperheight=900pt,margin=10pt]{geometry}
\usepackage{amsmath,amssymb}
\usepackage{array}
\usepackage{enumitem}
\usepackage{xcolor}
\usepackage{microtype}

\setlength{\parindent}{0pt}
\setlength{\parskip}{4pt}
\setlist[itemize]{leftmargin=12pt}
\setlist[enumerate]{leftmargin=14pt}

\sloppy
\allowdisplaybreaks
\emergencystretch=2em

\begin{document}

\section*{Uppgift 1}

\subsection*{Original question}
I ett stort land kandiderar Lena och Mats för att
bli president.
För enkelhetens skull antar vi att alla
röstberättigade medborgare vet vilken kandidat de
föredrar (\(L\) eller \(M\)) och att ingen vill lägga
ner sin röst.
Vi vill undersöka huruvida den politiska läggningen
hänger ihop med omständigheten att man är vegetarian
\((V)\) eller inte (ej \(V\)).
Följande sannolikheter är kända:

\[
P(L\cap V)=0.15,
\qquad
P(M\cap V)=0.05,
\qquad
P(L)=0.47.
\]

\begin{enumerate}[label=\alph*)]
\item Bestäm sannolikheterna
\(P(V)\), \(P(\text{ej }V)\) och \(P(L\mid V)\).
\hfill (1.5p)

\item Bestäm sannolikheten \(P(L\cup V)\).
\hfill (1p)

\item Är händelserna \(L\) och \(V\) oberoende?
\hfill (1p)

\item Man har frågat 10 personer (utan att utesluta
att samma person har frågats flera gånger) om deras
matvanor.
Vad är sannolikheten att fler än 5 har svarat att de
inte är vegetarianer?
\hfill (2.5p)
\end{enumerate}

\subsection*{Compiled notes from Yacoub + Obid}

\subsubsection*{What we are doing}
The notes split people by two things:
which candidate they vote for,
and whether they are vegetarian.

Then the notes use:
\begin{itemize}
\item conditional probability,
\item union formula,
\item independence check,
\item binomial distribution for 10 people.
\end{itemize}

\subsubsection*{Given information}
\[
P(L\cap V)=0.15.
\]

\[
P(M\cap V)=0.05.
\]

\[
P(L)=0.47.
\]

\[
L=\text{votes for Lena},
\qquad
M=\text{votes for Mats}.
\]

\[
V=\text{vegetarian}.
\]

\[
V^c=\text{not vegetarian}.
\]

\subsubsection*{Step-by-step solution}

\textbf{a) Find \(P(V)\)}

The notes say:
\[
P(V)
=
\text{all }V\text{ with }L
+
\text{all }V\text{ with }M.
\]

So:

\begin{align*}
P(V)
&=P(L\cap V)+P(M\cap V) \\[2mm]
&=0.15+0.05 \\[2mm]
&=0.20.
\end{align*}

So 20\% are vegetarian.

\textbf{a) Find \(P(\text{ej }V)\)}

\begin{align*}
P(V^c)
&=1-P(V) \\[2mm]
&=1-0.20 \\[2mm]
&=0.80.
\end{align*}

So 80\% are not vegetarian.

\textbf{a) Find \(P(L\mid V)\)}

The notes use the conditional formula:

\[
P(L\mid V)
=
\frac{P(L\cap V)}{P(V)}.
\]

Substitute:

\begin{align*}
P(L\mid V)
&=
\frac{0.15}{0.20} \\[2mm]
&=0.75.
\end{align*}

The notes explain this as:
knowing the person is vegetarian,
there is a 75\% chance they vote for \(L\).

\textbf{b) Find \(P(L\cup V)\)}

I ett stort land kandiderar Lena och Mats för att
bli president.
För enkelhetens skull antar vi att alla
röstberättigade medborgare vet vilken kandidat de
föredrar (\(L\) eller \(M\)) och att ingen vill lägga
ner sin röst.
Vi vill undersöka huruvida den politiska läggningen
hänger ihop med omständigheten att man är vegetarian
\((V)\) eller inte (ej \(V\)).
Följande sannolikheter är kända:

\[
P(L\cap V)=0.15,
\qquad
P(M\cap V)=0.05,
\qquad
P(L)=0.47.
\]

\begin{enumerate}[label=\alph*)]
\item Bestäm sannolikheterna
\(P(V)\), \(P(\text{ej }V)\) och \(P(L\mid V)\).
\hfill (1.5p)

\item Bestäm sannolikheten \(P(L\cup V)\).
\hfill (1p)

\item Är händelserna \(L\) och \(V\) oberoende?
\hfill (1p)

\item Man har frågat 10 personer (utan att utesluta
att samma person har frågats flera gånger) om deras
matvanor.
Vad är sannolikheten att fler än 5 har svarat att de
inte är vegetarianer?
\hfill (2.5p)
\end{enumerate}

The notes use the union formula:

\[
P(L\cup V)
=P(L)+P(V)-P(L\cap V).
\]

Substitute:

\begin{align*}
P(L\cup V)
&=0.47+0.20-0.15 \\[2mm]
&=0.52.
\end{align*}

\textbf{Explanation from the notes}

Intersection means:
people who voted \(L\) and are vegetarian.

Union means:
people who voted \(L\) or are vegetarian,
combined.

\textbf{c) Check independence}

I ett stort land kandiderar Lena och Mats för att
bli president.
För enkelhetens skull antar vi att alla
röstberättigade medborgare vet vilken kandidat de
föredrar (\(L\) eller \(M\)) och att ingen vill lägga
ner sin röst.
Vi vill undersöka huruvida den politiska läggningen
hänger ihop med omständigheten att man är vegetarian
\((V)\) eller inte (ej \(V\)).
Följande sannolikheter är kända:

\[
P(L\cap V)=0.15,
\qquad
P(M\cap V)=0.05,
\qquad
P(L)=0.47.
\]

\begin{enumerate}[label=\alph*)]
\item Bestäm sannolikheterna
\(P(V)\), \(P(\text{ej }V)\) och \(P(L\mid V)\).
\hfill (1.5p)

\item Bestäm sannolikheten \(P(L\cup V)\).
\hfill (1p)

\item Är händelserna \(L\) och \(V\) oberoende?
\hfill (1p)

\item Man har frågat 10 personer (utan att utesluta
att samma person har frågats flera gånger) om deras
matvanor.
Vad är sannolikheten att fler än 5 har svarat att de
inte är vegetarianer?
\hfill (2.5p)
\end{enumerate}

The notes check independence by comparing:

\[
P(L\cap V)
\quad \text{and} \quad
P(L)\cdot P(V).
\]

Use the values:

\begin{align*}
P(L)\cdot P(V)
&=0.47\cdot 0.20 \\[2mm]
&=0.094.
\end{align*}

But:

\[
P(L\cap V)=0.15.
\]

Compare:

\[
0.15\neq 0.094.
\]

Therefore:

\[
L \text{ and } V
\text{ are dependent.}
\]

\textbf{d) More than 5 are not vegetarian}

I ett stort land kandiderar Lena och Mats för att
bli president.
För enkelhetens skull antar vi att alla
röstberättigade medborgare vet vilken kandidat de
föredrar (\(L\) eller \(M\)) och att ingen vill lägga
ner sin röst.
Vi vill undersöka huruvida den politiska läggningen
hänger ihop med omständigheten att man är vegetarian
\((V)\) eller inte (ej \(V\)).
Följande sannolikheter är kända:

\[
P(L\cap V)=0.15,
\qquad
P(M\cap V)=0.05,
\qquad
P(L)=0.47.
\]

\begin{enumerate}[label=\alph*)]
\item Bestäm sannolikheterna
\(P(V)\), \(P(\text{ej }V)\) och \(P(L\mid V)\).
\hfill (1.5p)

\item Bestäm sannolikheten \(P(L\cup V)\).
\hfill (1p)

\item Är händelserna \(L\) och \(V\) oberoende?
\hfill (1p)

\item Man har frågat 10 personer (utan att utesluta
att samma person har frågats flera gånger) om deras
matvanor.
Vad är sannolikheten att fler än 5 har svarat att de
inte är vegetarianer?
\hfill (2.5p)
\end{enumerate}

The notes use:

\[
n=10.
\]

From part a):

\[
P(V^c)=0.80.
\]

Let:

\[
X=\text{number of not vegetarian persons}.
\]

Then:

\[
X\sim \operatorname{Bin}(10,0.80).
\]

We need:

\[
P(X>5).
\]

That means:

\[
X=6,7,8,9,10.
\]

The binomial formula is:

\[
P(X=k)
=
\binom{n}{k}
p^k(1-p)^{n-k}.
\]

For example, the notes calculate \(X=6\):

\begin{align*}
P(X=6)
&=
\binom{10}{6}
0.80^6(0.20)^4 \\[2mm]
&=210\cdot 0.26214\cdot 0.0016 \\[2mm]
&\approx 0.0881.
\end{align*}

The notes list:

\[
P(X=6)=0.0881.
\]

\[
P(X=7)=0.2013.
\]

\[
P(X=8)=0.3020.
\]

\[
P(X=9)=0.2684.
\]

\[
P(X=10)=0.1074.
\]

Add them:

\begin{align*}
P(X>5)
&=P(X=6)+P(X=7)+P(X=8) \\
&\quad +P(X=9)+P(X=10) \\[2mm]
&=0.0881+0.2013+0.3020 \\
&\quad +0.2684+0.1074 \\[2mm]
&=0.9672.
\end{align*}

So the notes conclude:

\[
P(X>5)=0.9672.
\]

\subsubsection*{Final answer}
\begin{enumerate}[label=\alph*)]
\item
\[
P(V)=0.20,
\qquad
P(\text{ej }V)=0.80.
\]

\[
P(L\mid V)=0.75.
\]

\item
\[
P(L\cup V)=0.52.
\]

\item
\[
L \text{ and } V
\text{ are dependent.}
\]

\item
\[
P(X>5)=0.9672.
\]

There is a 96.72\% chance that more than 5 out of
10 are not vegetarian.
\end{enumerate}

\subsubsection*{What to remember}
For conditional probability:

\[
P(A\mid B)
=
\frac{P(A\cap B)}{P(B)}.
\]

For a union:

\[
P(A\cup B)
=P(A)+P(B)-P(A\cap B).
\]

For independence:

\[
P(A\cap B)=P(A)P(B).
\]

For repeated yes/no trials:

\[
X\sim \operatorname{Bin}(n,p).
\]

\[
P(X=k)
=
\binom{n}{k}p^k(1-p)^{n-k}.
\]

\section*{Uppgift 2}

\subsection*{Original question}
Den stokastiska variabeln \(X\) har täthetsfunktionen

\[
f_X(x)=
\begin{cases}
x+1, & -1\le x<0,\\
-x+1, & 0\le x\le 1,\\
0, & \text{annars}.
\end{cases}
\]

\begin{enumerate}[label=\alph*)]
\item Verifiera att \(f_X\) uppfyller kraven på en
täthetsfunktion.
\hfill (1.5p)

\item Beräkna \(P(X\le 0.5)\).
\hfill (1p)

\item Bestäm väntevärde och varians för \(X\).
\hfill (2.5p)

\item Låt \(Y\) vara en stokastisk variabel som är
likformigt fördelad på intervallet \((-1,1)\) och som
är oberoende av \(X\).
Beräkna \(P(\max(X,Y)\le 0)\).
\hfill (2p)

\item Bestäm \(V(X+Y)\).
\hfill (1p)
\end{enumerate}

\subsection*{Compiled notes from Yacoub + Obid}

\subsubsection*{What we are doing}
The notes start by checking that the given function is
a valid probability density function.

For a valid density function, the notes use two checks:
\begin{itemize}
\item \(f(x)\ge 0\) for all \(x\),
\item the total probability is \(1\).
\end{itemize}

\subsubsection*{Given information}
\[
f_X(x)=
\begin{cases}
x+1, & -1\le x<0,\\
-x+1, & 0\le x\le 1,\\
0, & \text{annars}.
\end{cases}
\]

\[
X=\text{stokastisk variabel}.
\]

\subsubsection*{Step-by-step solution}

\textbf{a) Verify that \(f_X\) is a density}

The notes say that a density function must satisfy:

\[
f_X(x)\ge 0
\quad \text{for all }x.
\]

and:

\[
\int_{-\infty}^{\infty} f_X(x)\,dx=1.
\]

\textbf{Check 1: non-negativity}

On the interval:

\[
-1\le x<0
\]

we have:

\[
f_X(x)=x+1.
\]

The notes say the lowest value is at \(x=-1\):

\begin{align*}
f_X(-1)
&=-1+1 \\[2mm]
&=0.
\end{align*}

As \(x\) increases from \(-1\) to \(0\),
the function increases positively.
So this part is never negative.

On the interval:

\[
0\le x\le 1
\]

we have:

\[
f_X(x)=-x+1.
\]

The notes say the highest value is at \(x=0\):

\begin{align*}
f_X(0)
&=-0+1 \\[2mm]
&=1.
\end{align*}

At the end of the interval:

\begin{align*}
f_X(1)
&=-1+1 \\[2mm]
&=0.
\end{align*}

So this part is also never negative.

Outside these intervals:

\[
f_X(x)=0.
\]

Therefore:

\[
f_X(x)\ge 0
\quad \text{for all }x.
\]

\textbf{Check 2: total probability}

The notes then check that the total area is \(1\).

\begin{align*}
\int_{-\infty}^{\infty} f_X(x)\,dx
&=
\int_{-1}^{0}(x+1)\,dx \\
&\quad+
\int_{0}^{1}(-x+1)\,dx.
\end{align*}

First part:

\begin{align*}
\int_{-1}^{0}(x+1)\,dx
&=
\left[
\frac{x^2}{2}+x
\right]_{-1}^{0} \\[2mm]
&=0-
\left(
\frac{1}{2}-1
\right) \\[2mm]
&=\frac{1}{2}.
\end{align*}

Second part:

\begin{align*}
\int_{0}^{1}(-x+1)\,dx
&=
\left[
-\frac{x^2}{2}+x
\right]_{0}^{1} \\[2mm]
&=
-\frac{1}{2}+1 \\[2mm]
&=\frac{1}{2}.
\end{align*}

Add the two areas:

\begin{align*}
\frac{1}{2}+\frac{1}{2}
&=1.
\end{align*}

So the total area under the density is \(1\).

Therefore \(f_X\) is a valid density function.

\textbf{b) Calculate \(P(X\le 0.5)\)}

Den stokastiska variabeln \(X\) har täthetsfunktionen

\[
f_X(x)=
\begin{cases}
x+1, & -1\le x<0,\\
-x+1, & 0\le x\le 1,\\
0, & \text{annars}.
\end{cases}
\]

\begin{enumerate}[label=\alph*)]
\item Verifiera att \(f_X\) uppfyller kraven på en
täthetsfunktion.
\hfill (1.5p)

\item Beräkna \(P(X\le 0.5)\).
\hfill (1p)

\item Bestäm väntevärde och varians för \(X\).
\hfill (2.5p)

\item Låt \(Y\) vara en stokastisk variabel som är
likformigt fördelad på intervallet \((-1,1)\) och som
är oberoende av \(X\).
Beräkna \(P(\max(X,Y)\le 0)\).
\hfill (2p)

\item Bestäm \(V(X+Y)\).
\hfill (1p)
\end{enumerate}

The notes split the integral at \(0\), because the
formula for \(f_X\) changes there:

\begin{align*}
P(X\le 0.5)
&=
\int_{-1}^{0}(x+1)\,dx \\
&\quad+
\int_{0}^{0.5}(-x+1)\,dx.
\end{align*}

The first part is already known from part a):

\[
\int_{-1}^{0}(x+1)\,dx=\frac{1}{2}.
\]

Now calculate the second part:

\begin{align*}
\int_{0}^{0.5}(-x+1)\,dx
&=
\left[
-\frac{x^2}{2}+x
\right]_{0}^{0.5} \\[2mm]
&=
-\frac{0.5^2}{2}+0.5 \\[2mm]
&=
-0.125+0.5 \\[2mm]
&=0.375.
\end{align*}

Add the two parts:

\begin{align*}
P(X\le 0.5)
&=0.5+0.375 \\[2mm]
&=0.875.
\end{align*}

\textbf{c) Expected value and variance}

Den stokastiska variabeln \(X\) har täthetsfunktionen

\[
f_X(x)=
\begin{cases}
x+1, & -1\le x<0,\\
-x+1, & 0\le x\le 1,\\
0, & \text{annars}.
\end{cases}
\]

\begin{enumerate}[label=\alph*)]
\item Verifiera att \(f_X\) uppfyller kraven på en
täthetsfunktion.
\hfill (1.5p)

\item Beräkna \(P(X\le 0.5)\).
\hfill (1p)

\item Bestäm väntevärde och varians för \(X\).
\hfill (2.5p)

\item Låt \(Y\) vara en stokastisk variabel som är
likformigt fördelad på intervallet \((-1,1)\) och som
är oberoende av \(X\).
Beräkna \(P(\max(X,Y)\le 0)\).
\hfill (2p)

\item Bestäm \(V(X+Y)\).
\hfill (1p)
\end{enumerate}

The notes calculate:

\[
E[X]
=
\int x f_X(x)\,dx.
\]

Because \(f_X\) is piecewise:

\begin{align*}
E[X]
&=
\int_{-1}^{0}x(x+1)\,dx \\
&\quad+
\int_{0}^{1}x(-x+1)\,dx.
\end{align*}

Rewrite inside the integrals:

\begin{align*}
E[X]
&=
\int_{-1}^{0}(x^2+x)\,dx \\
&\quad+
\int_{0}^{1}(-x^2+x)\,dx.
\end{align*}

Calculate:

\begin{align*}
E[X]
&=
\left[
\frac{x^3}{3}
+\frac{x^2}{2}
\right]_{-1}^{0} \\
&\quad+
\left[
-\frac{x^3}{3}
+\frac{x^2}{2}
\right]_{0}^{1} \\[2mm]
&=
-\frac{1}{6}
+\frac{1}{6} \\[2mm]
&=0.
\end{align*}

So:

\[
E[X]=0.
\]

For the variance, the notes first find \(E[X^2]\):

\begin{align*}
E[X^2]
&=
\int_{-1}^{0}x^2(x+1)\,dx \\
&\quad+
\int_{0}^{1}x^2(-x+1)\,dx.
\end{align*}

Rewrite:

\begin{align*}
E[X^2]
&=
\int_{-1}^{0}(x^3+x^2)\,dx \\
&\quad+
\int_{0}^{1}(-x^3+x^2)\,dx.
\end{align*}

Calculate:

\begin{align*}
E[X^2]
&=
\left[
\frac{x^4}{4}
+\frac{x^3}{3}
\right]_{-1}^{0} \\
&\quad+
\left[
-\frac{x^4}{4}
+\frac{x^3}{3}
\right]_{0}^{1} \\[2mm]
&=
\frac{1}{12}
+\frac{1}{12} \\[2mm]
&=
\frac{1}{6}.
\end{align*}

The notes then use:

\[
\operatorname{Var}(X)
=
E[X^2]-(E[X])^2.
\]

So:

\begin{align*}
\operatorname{Var}(X)
&=
\frac{1}{6}-0^2 \\[2mm]
&=
\frac{1}{6}.
\end{align*}

The notes also write the standard deviation:

\begin{align*}
\sigma_X
&=
\sqrt{\frac{1}{6}} \\[2mm]
&=
\frac{\sqrt{6}}{6}.
\end{align*}

\textbf{d) Calculate \(P(\max(X,Y)\le 0)\)}

Den stokastiska variabeln \(X\) har täthetsfunktionen

\[
f_X(x)=
\begin{cases}
x+1, & -1\le x<0,\\
-x+1, & 0\le x\le 1,\\
0, & \text{annars}.
\end{cases}
\]

\begin{enumerate}[label=\alph*)]
\item Verifiera att \(f_X\) uppfyller kraven på en
täthetsfunktion.
\hfill (1.5p)

\item Beräkna \(P(X\le 0.5)\).
\hfill (1p)

\item Bestäm väntevärde och varians för \(X\).
\hfill (2.5p)

\item Låt \(Y\) vara en stokastisk variabel som är
likformigt fördelad på intervallet \((-1,1)\) och som
är oberoende av \(X\).
Beräkna \(P(\max(X,Y)\le 0)\).
\hfill (2p)

\item Bestäm \(V(X+Y)\).
\hfill (1p)
\end{enumerate}

The notes rewrite the event:

\[
\max(X,Y)\le 0
\]

means:

\[
X\le 0
\quad \text{and} \quad
Y\le 0.
\]

So:

\[
P(\max(X,Y)\le 0)
=
P(X\le 0,\ Y\le 0).
\]

The question says \(X\) and \(Y\) are independent.
The notes use:

\[
P(X\le 0,\ Y\le 0)
=
P(X\le 0)\cdot P(Y\le 0).
\]

From the earlier area calculation:

\[
P(X\le 0)=\frac{1}{2}.
\]

Since \(Y\) is uniformly distributed on \((-1,1)\),
exactly half of the interval is below \(0\):

\[
P(Y\le 0)=\frac{1}{2}.
\]

Therefore:

\begin{align*}
P(\max(X,Y)\le 0)
&=
\frac{1}{2}\cdot\frac{1}{2} \\[2mm]
&=
\frac{1}{4}.
\end{align*}

\textbf{e) Calculate \(V(X+Y)\)}

Den stokastiska variabeln \(X\) har täthetsfunktionen

\[
f_X(x)=
\begin{cases}
x+1, & -1\le x<0,\\
-x+1, & 0\le x\le 1,\\
0, & \text{annars}.
\end{cases}
\]

\begin{enumerate}[label=\alph*)]
\item Verifiera att \(f_X\) uppfyller kraven på en
täthetsfunktion.
\hfill (1.5p)

\item Beräkna \(P(X\le 0.5)\).
\hfill (1p)

\item Bestäm väntevärde och varians för \(X\).
\hfill (2.5p)

\item Låt \(Y\) vara en stokastisk variabel som är
likformigt fördelad på intervallet \((-1,1)\) och som
är oberoende av \(X\).
Beräkna \(P(\max(X,Y)\le 0)\).
\hfill (2p)

\item Bestäm \(V(X+Y)\).
\hfill (1p)
\end{enumerate}

The notes say that we already know:

\[
\operatorname{Var}(X)=\frac{1}{6}.
\]

Now find \(\operatorname{Var}(Y)\).
For a uniform distribution \(U(a,b)\), the notes use:

\[
\operatorname{Var}(Y)
=
\frac{(b-a)^2}{12}.
\]

Here:

\[
a=-1,
\qquad
b=1.
\]

So:

\begin{align*}
\operatorname{Var}(Y)
&=
\frac{(1-(-1))^2}{12} \\[2mm]
&=
\frac{2^2}{12} \\[2mm]
&=
\frac{4}{12} \\[2mm]
&=
\frac{1}{3}.
\end{align*}

The notes also warn that
\(P(Y\le 0)\) and \(\operatorname{Var}(Y)\)
are different concepts.

Since \(X\) and \(Y\) are independent:

\begin{align*}
\operatorname{Var}(X+Y)
&=
\operatorname{Var}(X)
+
\operatorname{Var}(Y) \\[2mm]
&=
\frac{1}{6}
+
\frac{1}{3} \\[2mm]
&=
\frac{1}{2}.
\end{align*}

\subsubsection*{Final answer}
\begin{enumerate}[label=\alph*)]
\item
\[
f_X
\text{ is a valid density function.}
\]

\item
\[
P(X\le 0.5)=0.875.
\]

\item
\[
E[X]=0,
\qquad
\operatorname{Var}(X)=\frac{1}{6}.
\]

\item
\[
P(\max(X,Y)\le 0)=\frac{1}{4}.
\]

\item
\[
V(X+Y)=\frac{1}{2}.
\]
\end{enumerate}

\subsubsection*{What to remember}
A density function must satisfy:

\[
f(x)\ge 0
\quad \text{for all }x.
\]

and:

\[
\int_{-\infty}^{\infty} f(x)\,dx=1.
\]

For a piecewise density, check each interval separately.

For independent variables:

\[
P(A\cap B)=P(A)P(B).
\]

For independent variables:

\[
\operatorname{Var}(X+Y)
=
\operatorname{Var}(X)
+
\operatorname{Var}(Y).
\]

For \(Y\sim U(a,b)\):

\[
\operatorname{Var}(Y)
=
\frac{(b-a)^2}{12}.
\]

\section*{Uppgift 3}

\subsection*{Original question}
Man kastar en tärning två gånger och noterar resultaten i tur och ordning.
Kasten kan anses som oberoende. Utfallsrummet kan då beskrivas att bestå av
utfallen \((i,j)\) där \(i\) är resultatet av första och \(j\) resultatet
av andra kastet, \(i,j=1,2,\ldots,6\).

\begin{enumerate}[label=\alph*)]
\item
Låt \(X\) vara den stokastiska variabeln med
\[
X(i,j)=|i-j|
\]
(t.ex. antar \(X\) värdet 3 om första kastet är en etta och andra kastet
är en fyra ty \(X(1,4)=|1-4|=3\)).
Bestäm sannolikheten \(P(X=4)\).

\item
Bestäm sannolikhetsfunktionen för \(X\) och beräkna \(E(X)\).

\item
I ett spel vinner spelaren 3 kronor om \(X\) minst är lika med 4 och
förlorar en krona annars. Lönar spelet sig för spelaren på lång sikt?
\end{enumerate}

\subsection*{Compiled notes from Yacoub + Obid}

\subsubsection*{What we are doing}
We throw one die two times.
The ordered result is written as \((i,j)\).
The random variable is the distance between the two throws:

\[
X=|i-j|.
\]

So \(X\) is discrete.
That means we count outcomes instead of integrating.

\subsubsection*{Given information}
\begin{itemize}
\item
First throw: \(i=1,2,3,4,5,6\).
\item
Second throw: \(j=1,2,3,4,5,6\).
\item
The throws are independent.
\item
Total number of ordered outcomes:
\[
6\cdot 6=36.
\]
\item
Possible values of \(X\):
\[
0,1,2,3,4,5.
\]
\item
Win condition in the game:
\[
X\ge 4.
\]
\end{itemize}

\subsubsection*{Step-by-step solution}

\textbf{a) Find \(P(X=4)\)}

We need all ordered pairs where the difference is 4.

\begin{align*}
(1,5)&:\ |1-5|=4,\\
(5,1)&:\ |5-1|=4,\\
(2,6)&:\ |2-6|=4,\\
(6,2)&:\ |6-2|=4.
\end{align*}

So there are 4 cases where \(X=4\).
There are 36 outcomes in total.

\begin{align*}
P(X=4)
&=
\frac{4}{36}\\[2mm]
&=
\frac{1}{9}.
\end{align*}

\textbf{b) Probability mass function for \(X\)}

Man kastar en tärning två gånger och noterar resultaten i tur och ordning.
Kasten kan anses som oberoende. Utfallsrummet kan då beskrivas att bestå av
utfallen \((i,j)\) där \(i\) är resultatet av första och \(j\) resultatet
av andra kastet, \(i,j=1,2,\ldots,6\).

\begin{enumerate}[label=\alph*)]
\item
Låt \(X\) vara den stokastiska variabeln med
\[
X(i,j)=|i-j|
\]
(t.ex. antar \(X\) värdet 3 om första kastet är en etta och andra kastet
är en fyra ty \(X(1,4)=|1-4|=3\)).
Bestäm sannolikheten \(P(X=4)\).

\item
Bestäm sannolikhetsfunktionen för \(X\) och beräkna \(E(X)\).

\item
I ett spel vinner spelaren 3 kronor om \(X\) minst är lika med 4 och
förlorar en krona annars. Lönar spelet sig för spelaren på lång sikt?
\end{enumerate}

The notes count how many pairs give each possible value of \(X\).

For \(X=0\), both dice are the same:

\[
(1,1),(2,2),(3,3),
(4,4),(5,5),(6,6).
\]

So:

\[
P(X=0)=\frac{6}{36}.
\]

For \(X=1\), the notes list 10 outcomes:

\[
(1,2),(2,1),(2,3),(3,2),(3,4),
\]
\[
(4,3),(5,4),(4,5),(5,6),(6,5).
\]

So:

\[
P(X=1)=\frac{10}{36}.
\]

The remaining counted outcomes are:

\[
P(X=2)=\frac{8}{36},
\]
\[
P(X=3)=\frac{6}{36},
\]
\[
P(X=4)=\frac{4}{36},
\]
\[
P(X=5)=\frac{2}{36}.
\]

So the probability mass function is:

\begin{align*}
p_X(0)&=\frac{6}{36},\\
p_X(1)&=\frac{10}{36},\\
p_X(2)&=\frac{8}{36},\\
p_X(3)&=\frac{6}{36},\\
p_X(4)&=\frac{4}{36},\\
p_X(5)&=\frac{2}{36}.
\end{align*}

\textbf{Calculate \(E(X)\)}

The notes use:

\[
E(X)=\sum x\cdot P(X=x).
\]

\begin{align*}
E(X)
&=
0\cdot\frac{6}{36}
+1\cdot\frac{10}{36}\\
&\quad
+2\cdot\frac{8}{36}
+3\cdot\frac{6}{36}\\
&\quad
+4\cdot\frac{4}{36}
+5\cdot\frac{2}{36}\\[2mm]
&=
\frac{0+10+16+18+16+10}{36}\\[2mm]
&=
\frac{70}{36}\\[2mm]
&=
\frac{35}{18}\\[2mm]
&\approx
1.94.
\end{align*}

\textbf{c) Long-run profit of the game}

Man kastar en tärning två gånger och noterar resultaten i tur och ordning.
Kasten kan anses som oberoende. Utfallsrummet kan då beskrivas att bestå av
utfallen \((i,j)\) där \(i\) är resultatet av första och \(j\) resultatet
av andra kastet, \(i,j=1,2,\ldots,6\).

\begin{enumerate}[label=\alph*)]
\item
Låt \(X\) vara den stokastiska variabeln med
\[
X(i,j)=|i-j|
\]
(t.ex. antar \(X\) värdet 3 om första kastet är en etta och andra kastet
är en fyra ty \(X(1,4)=|1-4|=3\)).
Bestäm sannolikheten \(P(X=4)\).

\item
Bestäm sannolikhetsfunktionen för \(X\) och beräkna \(E(X)\).

\item
I ett spel vinner spelaren 3 kronor om \(X\) minst är lika med 4 och
förlorar en krona annars. Lönar spelet sig för spelaren på lång sikt?
\end{enumerate}

The notes say:

\begin{itemize}
\item
Win scenario: \(X\ge 4\), win 3 kr.
\item
Lose scenario: \(X<4\), lose 1 kr.
\end{itemize}

First find the chance of winning:

\begin{align*}
P(X\ge 4)
&=
P(X=4)+P(X=5)\\[2mm]
&=
\frac{4}{36}
+
\frac{2}{36}\\[2mm]
&=
\frac{6}{36}\\[2mm]
&=
\frac{1}{6}.
\end{align*}

So the chance of losing is:

\begin{align*}
P(X<4)
&=
1-P(X\ge 4)\\[2mm]
&=
1-\frac{1}{6}\\[2mm]
&=
\frac{5}{6}.
\end{align*}

Now calculate expected profit per game.

\begin{align*}
E(W)
&=
(\text{win amount})P(\text{win})\\
&\quad+
(\text{loss amount})P(\text{loss})\\[2mm]
&=
3\cdot\frac{1}{6}
+
(-1)\cdot\frac{5}{6}\\[2mm]
&=
\frac{3}{6}
-
\frac{5}{6}\\[2mm]
&=
-\frac{2}{6}\\[2mm]
&=
-\frac{1}{3}.
\end{align*}

Because the expected profit is negative, the player will not gain money
in the long run.

\subsubsection*{Final answer}
\begin{enumerate}[label=\alph*)]
\item
\[
P(X=4)=\frac{1}{9}.
\]

\item
\[
p_X(0)=\frac{6}{36},
\quad
p_X(1)=\frac{10}{36},
\]
\[
p_X(2)=\frac{8}{36},
\quad
p_X(3)=\frac{6}{36},
\]
\[
p_X(4)=\frac{4}{36},
\quad
p_X(5)=\frac{2}{36}.
\]

\[
E(X)=\frac{35}{18}\approx 1.94.
\]

\item
\[
E(W)=-\frac{1}{3}.
\]

The game does not pay off for the player in the long run.
\end{enumerate}

\subsubsection*{What to remember}
When a random variable is discrete, use a probability mass function:

\[
p_X(x)=P(X=x).
\]

For dice problems, count outcomes:

\[
P(\text{event})
=
\frac{\text{number of good outcomes}}
{\text{number of all outcomes}}.
\]

Expected value for a discrete variable is:

\[
E(X)=\sum x\cdot P(X=x).
\]

For a game, check if it is profitable by calculating expected profit.
If the expected profit is negative, it is bad for the player in the long run.

\section*{Uppgift 4}

\subsection*{Original question}

Vi studerar mätningar som är fördelade enligt
\(N(\mu,0.5)\) där \(\mu\) är det exakta värdet
och \(\sigma = 0.5\) återspeglar mätinstrumentets
precision.

\begin{enumerate}[label=\alph*)]
\item Beräkna \(P(X \le 10.375)\) för en mätning
\(X \in N(\mu,0.5)\) där \(\mu = 10\). \((1p)\)

\item Antag att vi har ett stickprov
\(x_1,\ldots,x_5\) av 5 observerade mätningar som
antas komma från \(N(\mu,0.5)\) med samma \(\mu\).
Utifrån stickprovet har vi beräknat \(\bar{x}=9.6\),
bestäm ett (två-sidigt) 95\%-konfidensintervall
för \(\mu\). \((2p)\)

\item Testa hypotesen \(\mu = 10\) på
5\% signifikansnivå. \((1p)\)

\item Testa hypotesen \(\mu \le 10\) på
5\% signifikansnivå. \((2p)\)
\end{enumerate}

\subsection*{Compiled notes from Yacoub + Obid}
% UNCERTAIN: The supplied note images are not visibly labelled by author,
% so Yacoub's and Obid's notes are compiled together as the provided notes.

\subsubsection*{What we are doing}

We use the normal distribution with known
\(\sigma=0.5\). First we standardize to a
\(Z\)-score. Then we use the \(Z\)-table.

For the interval and tests, we use the sample data
from part b):
\[
\bar{x}=9.6,\qquad n=5,\qquad \sigma=0.5.
\]

\subsubsection*{Given information}

\[
X \sim N(\mu,0.5)
\]

\[
\sigma = 0.5
\]

For part a):
\[
\mu = 10,
\qquad
X \le 10.375.
\]

For parts b)--d):
\[
n=5,
\qquad
\bar{x}=9.6.
\]

\subsubsection*{Step-by-step solution}

\textbf{a) Find \(P(X \le 10.375)\)}

The notes standardize \(X\) to a \(Z\)-score:
\[
Z=\frac{X-\mu}{\sigma}.
\]

Substitute the numbers:
\begin{align*}
z
&=
\frac{10.375-10}{0.5}
\\[2mm]
&=
\frac{0.375}{0.5}
\\[2mm]
&=0.75.
\end{align*}

So
\begin{align*}
P(X \le 10.375)
&=
P(Z \le 0.75).
\end{align*}

From the \(Z\)-table, read row \(0.7\) and
column \(0.05\):
\[
P(Z \le 0.75)=0.7734.
\]

\textbf{b) 95\% confidence interval for \(\mu\)}

Vi studerar mätningar som är fördelade enligt
\(N(\mu,0.5)\) där \(\mu\) är det exakta värdet
och \(\sigma = 0.5\) återspeglar mätinstrumentets
precision.

\begin{enumerate}[label=\alph*)]
\item Beräkna \(P(X \le 10.375)\) för en mätning
\(X \in N(\mu,0.5)\) där \(\mu = 10\). \((1p)\)

\item Antag att vi har ett stickprov
\(x_1,\ldots,x_5\) av 5 observerade mätningar som
antas komma från \(N(\mu,0.5)\) med samma \(\mu\).
Utifrån stickprovet har vi beräknat \(\bar{x}=9.6\),
bestäm ett (två-sidigt) 95\%-konfidensintervall
för \(\mu\). \((2p)\)

\item Testa hypotesen \(\mu = 10\) på
5\% signifikansnivå. \((1p)\)

\item Testa hypotesen \(\mu \le 10\) på
5\% signifikansnivå. \((2p)\)
\end{enumerate}

From the question and notes:
\[
n=5,\qquad \sigma=0.5,\qquad \bar{x}=9.6.
\]

It is a two-sided 95\% confidence interval.
So 5\% is outside the interval, split into
2.5\% in each tail.

The notes use:
\[
z_{0.025}\approx 1.96.
\]

First calculate the standard error:
\begin{align*}
SE
&=
\frac{\sigma}{\sqrt{n}}
\\[2mm]
&=
\frac{0.5}{\sqrt{5}}
\\[2mm]
&\approx 0.2236.
\end{align*}

Then calculate the margin of error:
\begin{align*}
ME
&=
z\cdot SE
\\[2mm]
&=
1.96\cdot 0.2236
\\[2mm]
&\approx 0.436.
\end{align*}

The interval is:
\[
\bar{x}\pm ME.
\]

Substitute:
\begin{align*}
9.6 \pm 0.436
&=
\bigl[
9.6-0.436,\,
9.6+0.436
\bigr]
\\[2mm]
&=
[9.164,\,10.036].
\end{align*}

\textbf{English clarification:}
The note image writes the final interval as
\([9.112,\,10.038]\), but the written calculation
\(9.6\pm 0.436\) gives \([9.164,\,10.036]\).
The method from the notes is kept here.
% UNCERTAIN: The handwritten final lower endpoint looks like 9.112,
% but 9.6 - 0.436 = 9.164.

\textbf{c) Test \(\mu=10\) at 5\% significance}

Vi studerar mätningar som är fördelade enligt
\(N(\mu,0.5)\) där \(\mu\) är det exakta värdet
och \(\sigma = 0.5\) återspeglar mätinstrumentets
precision.

\begin{enumerate}[label=\alph*)]
\item Beräkna \(P(X \le 10.375)\) för en mätning
\(X \in N(\mu,0.5)\) där \(\mu = 10\). \((1p)\)

\item Antag att vi har ett stickprov
\(x_1,\ldots,x_5\) av 5 observerade mätningar som
antas komma från \(N(\mu,0.5)\) med samma \(\mu\).
Utifrån stickprovet har vi beräknat \(\bar{x}=9.6\),
bestäm ett (två-sidigt) 95\%-konfidensintervall
för \(\mu\). \((2p)\)

\item Testa hypotesen \(\mu = 10\) på
5\% signifikansnivå. \((1p)\)

\item Testa hypotesen \(\mu \le 10\) på
5\% signifikansnivå. \((2p)\)
\end{enumerate}

The notes set up the two-sided hypotheses:
\[
H_0:\mu=10
\]
\[
H_1:\mu\ne 10.
\]

Use the sample data from part b).
The test statistic is:
\[
Z=
\frac{\bar{x}-\mu_0}{\sigma/\sqrt{n}}.
\]

Substitute:
\begin{align*}
Z
&=
\frac{9.6-10}{0.5/\sqrt{5}}
\\[2mm]
&=
\frac{-0.4}{0.2236}
\\[2mm]
&\approx -1.789.
\end{align*}

For a two-sided 5\% test, the notes use the
critical values:
\[
-1.96
\quad\text{and}\quad
1.96.
\]

The test statistic is inside the interval:
\[
-1.96 < -1.789 < 1.96.
\]

So the notes conclude that \(\mu=10\) is accepted.

\textbf{English clarification:}
In hypothesis-test language, this means we do not
reject \(H_0\) at the 5\% level.

\textbf{d) Test \(\mu \le 10\) at 5\% significance}

Vi studerar mätningar som är fördelade enligt
\(N(\mu,0.5)\) där \(\mu\) är det exakta värdet
och \(\sigma = 0.5\) återspeglar mätinstrumentets
precision.

\begin{enumerate}[label=\alph*)]
\item Beräkna \(P(X \le 10.375)\) för en mätning
\(X \in N(\mu,0.5)\) där \(\mu = 10\). \((1p)\)

\item Antag att vi har ett stickprov
\(x_1,\ldots,x_5\) av 5 observerade mätningar som
antas komma från \(N(\mu,0.5)\) med samma \(\mu\).
Utifrån stickprovet har vi beräknat \(\bar{x}=9.6\),
bestäm ett (två-sidigt) 95\%-konfidensintervall
för \(\mu\). \((2p)\)

\item Testa hypotesen \(\mu = 10\) på
5\% signifikansnivå. \((1p)\)

\item Testa hypotesen \(\mu \le 10\) på
5\% signifikansnivå. \((2p)\)
\end{enumerate}

The notes say this is the same as before, but now
the hypothesis is one-sided:
\[
H_0:\mu\le 10
\]
\[
H_1:\mu>10.
\]

This is a right-tailed test.

From part c), we still have:
\[
Z\approx -1.789.
\]

For a right-tailed 5\% test, the notes use:
\[
z_{0.05}=1.645.
\]

Reject \(H_0\) if the test statistic is bigger
than \(1.645\).

But:
\[
-1.789 < 1.645.
\]

So the notes conclude that \(H_0\) is accepted.

\textbf{English clarification:}
In hypothesis-test language, this means we do not
reject \(H_0:\mu\le 10\) at the 5\% level.

\subsubsection*{Final answer}

\begin{enumerate}[label=\alph*)]
\item
\[
P(X\le 10.375)=0.7734.
\]

\item
Using the written method:
\[
9.6\pm 0.436
= [9.164,\,10.036].
\]
% UNCERTAIN: Handwritten final answer appears to show
% [9.112, 10.038], but this does not match 9.6 +/- 0.436.

\item
\[
Z\approx -1.789.
\]
Since \(-1.789\) lies between \(-1.96\) and \(1.96\),
we do not reject \(H_0:\mu=10\).

\item
\[
Z\approx -1.789.
\]
Since \(-1.789<1.645\), we do not reject
\(H_0:\mu\le 10\).
\end{enumerate}

\subsubsection*{What to remember}

When \(\sigma\) is known, the notes use the
\(Z\)-formula:
\[
Z=
\frac{\bar{x}-\mu_0}{\sigma/\sqrt{n}}.
\]

For a 95\% two-sided interval:
\[
\bar{x}
\pm
1.96\cdot
\frac{\sigma}{\sqrt{n}}.
\]

For a right-tailed 5\% test, compare the test
statistic with:
\[
1.645.
\]

\section*{Uppgift 5}

\subsection*{Original question}

Antag att vi har ett stickprov
\[
0<x_1,\ldots,x_n<1
\]
från en kontinuerlig fördelning med täthetsfunktionen

\[
f(x)=
\begin{cases}
\theta x^{\theta-1},
& 0<x<1,\\
0,
& \text{annars},
\end{cases}
\]

där \(\theta>0\).

\begin{enumerate}[label=\alph*)]
\item Ange likelihoodfunktionen \(L(\theta)\).
\hfill \((1.5p)\)

\item Bestäm ML-skattningen
\(\theta^*_{\text{obs}}\) av \(\theta\).
\hfill \((4.5p)\)
\end{enumerate}

\textit{Tipps:} I stället för att maximera
\(L(\theta)\) kan man maximera \(\ln L(\theta)\).
Kan du motivera detta?

\subsection*{Compiled notes from Yacoub + Obid}
% UNCERTAIN: The supplied note images are not visibly labelled by author,
% so Yacoub's and Obid's notes are compiled together as the provided notes.

\subsubsection*{What we are doing}

We have a sample
\[
0<x_1,\ldots,x_n<1.
\]

The notes first build the likelihood by multiplying
the density values:
\[
L(\theta)=\prod_{i=1}^{n} f(x_i).
\]

Then the notes use \(\ln(L(\theta))\), because the
logarithm turns products and powers into sums.
That makes the derivative easier.

\subsubsection*{Given information}

\[
f(x)=
\begin{cases}
\theta x^{\theta-1},
& 0<x<1,\\
0,
& \text{annars}.
\end{cases}
\]

\[
\theta>0.
\]

\[
0<x_i<1
\quad
\text{for all }i.
\]

\subsubsection*{Step-by-step solution}

\textbf{a) Likelihood function}

The notes write:
\[
L(\theta)
=
\prod_{i=1}^{n} f(x_i).
\]

Since every observation satisfies \(0<x_i<1\),
we use:
\[
f(x_i)=\theta x_i^{\theta-1}.
\]

So:
\begin{align*}
L(\theta)
&=
\prod_{i=1}^{n}
\left(
\theta x_i^{\theta-1}
\right)
\\[2mm]
&=
\prod_{i=1}^{n}\theta
\cdot
\prod_{i=1}^{n}
x_i^{\theta-1}
\\[2mm]
&=
\theta^n
\prod_{i=1}^{n}
x_i^{\theta-1}.
\end{align*}

Therefore:
\[
L(\theta)
=
\theta^n
\prod_{i=1}^{n}
x_i^{\theta-1},
\qquad
\theta>0.
\]

\textbf{b) ML estimate}

Antag att vi har ett stickprov
\[
0<x_1,\ldots,x_n<1
\]
från en kontinuerlig fördelning med täthetsfunktionen

\[
f(x)=
\begin{cases}
\theta x^{\theta-1},
& 0<x<1,\\
0,
& \text{annars},
\end{cases}
\]

där \(\theta>0\).

\begin{enumerate}[label=\alph*)]
\item Ange likelihoodfunktionen \(L(\theta)\).
\hfill \((1.5p)\)

\item Bestäm ML-skattningen
\(\theta^*_{\text{obs}}\) av \(\theta\).
\hfill \((4.5p)\)
\end{enumerate}

\textit{Tipps:} I stället för att maximera
\(L(\theta)\) kan man maximera \(\ln L(\theta)\).
Kan du motivera detta?

The notes say that \(\ln\) helps remove the annoying
exponents, so we use:
\[
\ln(L(\theta)).
\]

Start from the likelihood:
\[
L(\theta)
=
\theta^n
\prod_{i=1}^{n}
x_i^{\theta-1}.
\]

Take logarithms:
\begin{align*}
\ln(L(\theta))
&=
\ln\left(
\theta^n
\prod_{i=1}^{n}
x_i^{\theta-1}
\right)
\\[2mm]
&=
\ln(\theta^n)
+
\ln\left(
\prod_{i=1}^{n}
x_i^{\theta-1}
\right).
\end{align*}

Use the log rules from the notes:
\[
\ln(AB)=\ln(A)+\ln(B),
\]
\[
\ln(\theta^n)=n\ln(\theta),
\]
\[
\ln\left(\prod f(x_i)\right)
=
\sum \ln(f(x_i)).
\]

Then:
\begin{align*}
\ln(L(\theta))
&=
n\ln(\theta)
+
\sum_{i=1}^{n}
\ln\left(
x_i^{\theta-1}
\right)
\\[2mm]
&=
n\ln(\theta)
+
\sum_{i=1}^{n}
(\theta-1)\ln(x_i)
\\[2mm]
&=
n\ln(\theta)
+
(\theta-1)
\sum_{i=1}^{n}
\ln(x_i).
\end{align*}

Now differentiate with respect to \(\theta\):
\begin{align*}
\frac{d}{d\theta}
\ln(L(\theta))
&=
\frac{d}{d\theta}
\left[
n\ln(\theta)
+
(\theta-1)
\sum_{i=1}^{n}
\ln(x_i)
\right]
\\[2mm]
&=
\frac{n}{\theta}
+
1\cdot
\sum_{i=1}^{n}
\ln(x_i)
\\[2mm]
&=
\frac{n}{\theta}
+
\sum_{i=1}^{n}
\ln(x_i).
\end{align*}

Set the derivative equal to zero:
\begin{align*}
\frac{n}{\theta}
+
\sum_{i=1}^{n}
\ln(x_i)
&=0.
\end{align*}

Move the sum to the other side:
\begin{align*}
\frac{n}{\theta}
&=
-
\sum_{i=1}^{n}
\ln(x_i).
\end{align*}

Solve for \(\theta\):
\begin{align*}
\theta^*_{\text{obs}}
&=
\frac{n}
{
-
\sum_{i=1}^{n}
\ln(x_i)
}.
\end{align*}

Equivalently:
\[
\theta^*_{\text{obs}}
=
-
\frac{n}
{
\sum_{i=1}^{n}
\ln(x_i)
}.
\]

\textbf{English clarification:}
Since \(0<x_i<1\), we have \(\ln(x_i)<0\).
Therefore
\[
-
\sum_{i=1}^{n}\ln(x_i)
\]
is positive, so the estimate is allowed because
\(\theta>0\).

\textbf{Why we may maximize \(\ln L(\theta)\)}

The logarithm is strictly increasing. That means it
keeps the same order:
\[
L(\theta_1)>L(\theta_2)
\quad \Longleftrightarrow \quad
\ln L(\theta_1)>\ln L(\theta_2).
\]

So the \(\theta\) that maximizes \(L(\theta)\) also
maximizes \(\ln L(\theta)\).

\textbf{English clarification:}
To check that the critical point is a maximum:
\begin{align*}
\frac{d^2}{d\theta^2}
\ln(L(\theta))
&=
-
\frac{n}{\theta^2}
\\[2mm]
&<0.
\end{align*}

So the log-likelihood is concave in \(\theta\), and
the critical point gives the maximum.

\subsubsection*{Final answer}

\begin{enumerate}[label=\alph*)]
\item
\[
L(\theta)
=
\theta^n
\prod_{i=1}^{n}
x_i^{\theta-1},
\qquad
\theta>0.
\]

\item
\[
\theta^*_{\text{obs}}
=
\frac{n}
{
-
\sum_{i=1}^{n}
\ln(x_i)
}
=
-
\frac{n}
{
\sum_{i=1}^{n}
\ln(x_i)
}.
\]
\end{enumerate}

\subsubsection*{What to remember}

Likelihood for independent observations:
\[
L(\theta)
=
\prod_{i=1}^{n} f(x_i).
\]

The logarithm is useful because:
\[
\ln(AB)=\ln(A)+\ln(B),
\]
\[
\ln(a^b)=b\ln(a).
\]

For ML estimation:
\begin{enumerate}
\item Write \(L(\theta)\).
\item Take \(\ln(L(\theta))\).
\item Differentiate with respect to \(\theta\).
\item Set the derivative equal to zero.
\item Solve for \(\theta\).
\end{enumerate}

\section*{Uppgift 6}

\subsection*{Original question}

En informatiker har implementerat en
slumptalsgenerator avsedd att generera en diskret
likformig sannolikhetsfördelning på de hela talen
\[
\{1,2,3,4,5,6\}.
\]
I en testkörning skapade informatikern 1000 slumptal
med hjälp av slumptalsgeneratorn erhöll följande
resultat:

\[
\begin{array}{c|c}
\text{Slumptal} & \text{Frekvens}\\
\hline
1 & 166\\
2 & 168\\
3 & 167\\
4 & 184\\
5 & 158\\
6 & 157
\end{array}
\]

\begin{enumerate}[label=\alph*)]
\item Formulera en rimlig nollhypotes för att avgöra
om slumpgeneratorn fungerar. \((1.5p)\)

\item Genomför ett test av nollhypotesen på
5\% signifikansnivå. \((2.5p)\)
\end{enumerate}

\subsection*{Compiled notes from Yacoub + Obid}
% UNCERTAIN: The supplied note images are not visibly labelled by author,
% so Yacoub's and Obid's notes are compiled together as the provided notes.

\subsubsection*{What we are doing}

We test if the random number generator is fair.
If it is fair, each number \(1,2,3,4,5,6\)
should have the same chance.

The notes use a chi-square goodness-of-fit test.
We compare:
\[
\text{observed frequencies}
\]
with:
\[
\text{expected frequencies}.
\]

\subsubsection*{Given information}

The generator should give numbers:
\[
1,2,3,4,5,6.
\]

The test was run:
\[
n=1000
\]
times.

Observed frequencies:
\[
O_1=166,\quad O_2=168,\quad O_3=167,
\]
\[
O_4=184,\quad O_5=158,\quad O_6=157.
\]

Significance level:
\[
\alpha=0.05.
\]

\subsubsection*{Step-by-step solution}

\textbf{a) Formulate the hypotheses}

The notes say that if the generator is working as
intended, all numbers should have the same chance.

So the null hypothesis is:
\[
H_0:
p(1)=p(2)=p(3)=p(4)=p(5)=p(6)
=\frac{1}{6}.
\]

The alternative hypothesis is the opposite:
\[
H_1:
\text{the probabilities are not all equal.}
\]

\textbf{b) Chi-square test at 5\% level}

En informatiker har implementerat en
slumptalsgenerator avsedd att generera en diskret
likformig sannolikhetsfördelning på de hela talen
\[
\{1,2,3,4,5,6\}.
\]
I en testkörning skapade informatikern 1000 slumptal
med hjälp av slumptalsgeneratorn erhöll följande
resultat:

\[
\begin{array}{c|c}
\text{Slumptal} & \text{Frekvens}\\
\hline
1 & 166\\
2 & 168\\
3 & 167\\
4 & 184\\
5 & 158\\
6 & 157
\end{array}
\]

\begin{enumerate}[label=\alph*)]
\item Formulera en rimlig nollhypotes för att avgöra
om slumpgeneratorn fungerar. \((1.5p)\)

\item Genomför ett test av nollhypotesen på
5\% signifikansnivå. \((2.5p)\)
\end{enumerate}

If the generator is fair, each number should appear
\(\frac{1}{6}\) of the time out of 1000 trials.

So the expected frequency for each number is:
\begin{align*}
E_i
&=
n\cdot p
\\[2mm]
&=
1000\cdot \frac{1}{6}
\\[2mm]
&\approx 166.67.
\end{align*}

The notes use the chi-square test statistic:
\[
\chi^2
=
\sum
\frac{(O-E)^2}{E}.
\]

Here \(O\) is the observed frequency from the table,
and \(E\) is the expected frequency.

Now compare each observed frequency with \(166.67\).

For number \(1\):
\begin{align*}
\frac{(166-166.67)^2}{166.67}
&\approx 0.003.
\end{align*}

For number \(2\):
\begin{align*}
\frac{(168-166.67)^2}{166.67}
&\approx 0.011.
\end{align*}

For number \(3\):
\begin{align*}
\frac{(167-166.67)^2}{166.67}
&\approx 0.001.
\end{align*}

For number \(4\):
\begin{align*}
\frac{(184-166.67)^2}{166.67}
&\approx 1.803.
\end{align*}

For number \(5\):
\begin{align*}
\frac{(158-166.67)^2}{166.67}
&\approx 0.451.
\end{align*}

For number \(6\):
\begin{align*}
\frac{(157-166.67)^2}{166.67}
&\approx 0.561.
\end{align*}

Add all terms:
\begin{align*}
\chi^2
&\approx
0.003+0.011+0.001
\\
&\quad
+1.803+0.451+0.561
\\[2mm]
&\approx 2.83.
\end{align*}

\textbf{Degrees of freedom}

The notes use:
\[
\text{df}
=
\text{number of categories}-1.
\]

There are 6 categories, so:
\begin{align*}
\text{df}
&=
6-1
\\[2mm]
&=5.
\end{align*}

\textbf{Critical value}

From the chi-square table, for:
\[
\alpha=0.05,
\qquad
\text{df}=5,
\]
the notes use:
\[
\chi^2_{0.05,5}=11.07.
\]

Compare:
\[
2.83<11.07.
\]

The test statistic is less than the rejection limit.
So the notes conclude that the original hypothesis is
correct and the generator is fair.

\textbf{English clarification:}
In hypothesis-test language, this means we do not
reject \(H_0\) at the 5\% significance level.

\subsubsection*{Final answer}

\begin{enumerate}[label=\alph*)]
\item
\[
H_0:
p(1)=p(2)=p(3)=p(4)=p(5)=p(6)
=\frac{1}{6}.
\]
\[
H_1:
\text{the probabilities are not all equal.}
\]

\item
\[
E_i\approx 166.67
\]
for each number.

\[
\chi^2\approx 2.83.
\]

\[
\text{df}=5,
\qquad
\chi^2_{0.05,5}=11.07.
\]

Since:
\[
2.83<11.07,
\]
we do not reject \(H_0\). The notes conclude that
the random number generator is fair.
\end{enumerate}

\subsubsection*{What to remember}

For a chi-square goodness-of-fit test:
\[
\chi^2
=
\sum
\frac{(O-E)^2}{E}.
\]

If there are \(k\) categories and no parameters are
estimated from the data:
\[
\text{df}=k-1.
\]

If:
\[
\chi^2_{\text{obs}}
<
\chi^2_{\text{critical}},
\]
then we do not reject \(H_0\).

\end{document}
